
\begin{frame}[fragile]{Deux exemples de schémas et de bases normalisés}

Cette session présente deux exemples de schémas relationnels normalisés. 

\vfill

Les notions
illustrées sont:

\begin{redbullet}
  \item Les clés primaires
  \item Les clés étrangères et les associations entre tables
  \item Le nommage des relations et des attributs
\end{redbullet}

\vfill
\begin{blocgauche}
\alert{Ces diapositives correspondent au support en ligne disponible sur
le site \url{http://sql.bdpedia.fr/}}
\end{blocgauche}

\end{frame}


\begin{frame}{La base des voyageurs}

Le schéma est résumé ci-dessous. 
\vfill

\begin{redbullet}
  \item Voyageur (\textbf{idVoyageur}, nom, prénom, ville, région)
    \item  Séjour (\textbf{idSéjour}, \textit{idVoyageur}, \textit{codeLogement}, début, fin)
    \item Logement (\textbf{code}, nom, capacité, type, lieu)
    \item Activité (\textbf{\textit{codeLogement}, codeActivité}, description)
\end{redbullet}

\vfill

En \textbf{gras}, les clés primaires, en
\textit{italiques} les clés étrangères.

\begin{blocgauche}
\begin{remblock}{Nommage}
Totalement libre. Il est préférable d'utiliser des conventions homogènes.
\end{remblock}
\end{blocgauche}

\end{frame}

\begin{frame}[fragile]{Table des voyageurs et des logements}

\begin{tabular}{c|c|c|c|c}
\textbf{idVoyageur} & \textbf{nom} & \textbf{prénom} & \textbf{ville} & \textbf{région}\\ \hline 
10 & Fogg & Phileas & Ajaccio & Corse\\ \hline 
20 & Bouvier & Nicolas & Aurillac & Auvergne\\ \hline 
30 & David-Néel & Alexandra & Lhassa & Tibet\\ \hline 
40 & Stevenson & Robert Louis & Vannes & Bretagne\\ \hline 
\end{tabular}

\vfill

\begin{tabular}{c|c|c|c|c}
\textbf{code} & \textbf{nom} & \textbf{capacité} & \textbf{type} & \textbf{lieu}\\ \hline 
ca & Causses & 45 & Auberge & Cévennes\\ \hline 
ge & Génépi & 134 & Hôtel & Alpes\\ \hline 
pi & U Pinzutu & 10 & Gîte & Corse\\ \hline 
ta & Tabriz & 34 & Hôtel & Bretagne\\ \hline 
\end{tabular}
\end{frame}


\begin{frame}[fragile]{Table des séjours}

\begin{center}
\begin{tabular}{c|c|c|c|c}
\textbf{idSéjour} & \textbf{\textit{idVoyageur}} & \textbf{\textit{codeLogement}} & \textbf{début} & \textbf{fin}\\ \hline 
1 & 10 & pi & 20 & 20\\ \hline 
2 & 20 & ta & 21 & 22\\ \hline 
3 & 30 & ge & 2 & 3\\ \hline 
4 & 20 & pi & 19 & 23\\ \hline 
5 & 20 & ge & 22 & 24\\ \hline 
6 & 10 & pi & 10 & 12\\ \hline 
7 & 30 & ca & 13 & 18\\ \hline 
\end{tabular}
\end{center}

\vfill

\begin{blocgauche}
\begin{remblock}{Rôle des clés étrangères}
Chaque nuplet référence \alert{un} voyageur et \alert{un} logement.
Mais un voyageur ou un logement peut être
référencé plusieurs fois.
\end{remblock}
\end{blocgauche}

\end{frame}



\begin{frame}[fragile]{Table des activités}
\begin{center}
\begin{tabular}{c|c|c}
\textbf{\textit{codeLogement}} & \textbf{codeActivité} & \textbf{description}\\ \hline 
ca & Randonnée & Sorties d’une journée en groupe\\ \hline 
ge & Piscine & Nage loisir non encadrée\\ \hline 
ge & Ski & Sur piste uniquement\\ \hline 
pi & Plongée & Baptèmes et préparation des brevets\\ \hline 
pi & Voile & Pratique du dériveur et du catamaran\\ \hline 
\end{tabular}
\end{center}


\vfill

\begin{blocgauche}
\begin{remblock}{Clé primaire et étrangère }
Une clé étrangère peut être une partie d'une clé primaire.
\end{remblock}
\end{blocgauche}
\end{frame}

%%%%%%ù

%%%%%%%%%%%%
\begin{frame}{La base des films}

Clés primaires en \textbf{gras}, clés étrangères en \textit{italiques}.

\vfill
\begin{redbullet}
  \item Film (\textbf{idFilm}, titre, année, genre, résumé, \textit{idRéalisateur}, \textit{codePays})
  \item Artiste (\textbf{idArtiste}, nom, prénom, annéeNaissance)
    \item Rôle (\textbf{\textit{idFilm}, \textit{idActeur}}, nomRôle)
   \item Internaute (\textbf{email}, nom, prénom, région)
     \item Notation (\textbf{\textit{email}, \textit{idFilm}}, note)
   \item Pays (\textbf{code}, nom, langue)
\end{redbullet}


\vfill
\begin{blocgauche}
Rappel: le nommage des attributs est libre.
\end{blocgauche}
\end{frame}

\begin{frame}{Petit exemple}


\begin{small}
\begin{center}
\begin{tabular}{l|l|c|l|c|c}
\textbf{idFilm} & titre & année  & genre & \textit{idMES} & \textit{codePays}\\
\hline
20 & Impitoyable & 1992 & Western & 100 & USA \\
21 & Ennemi d'état & 1998 & Action & 102 & USA\\
\hline
\end{tabular}

\centerline{La table \tble{Film}}
\end{center}
\begin{center}
\begin{minipage}[b]{6cm}
\begin{tabular}{l|l|l|c}
\textbf{idArtiste} &  nom & prénom &  annéeNaiss \\
\hline
100 & Eastwood & Clint  & 1930 \\
101 & Hackman & Gene  & 1930 \\
102 & Scott & Tony  & 1930 \\
103 & Smith  & Will  & 1968 \\
\hline
\end{tabular}
\centerline{La table \tble{Artiste}}
\end{minipage} \hspace*{2cm} \
\begin{minipage}[b]{5cm}
\begin{tabular}{l|c|l}
\textbf{\textit{idFilm}} & \textbf{\textit{idActeur}}  & rôle \\
\hline
20  & 100  & William Munny\\
20 & 101  & Little Bill\\
21 & 101  & Bril\\
21 & 103  & Robert Dean\\
\hline
\end{tabular}

\centerline{La table \tble{Rôle}}
\end{minipage}
\end{center}
\end{small}
\end{frame}

\begin{frame}{À retenir}

Comprendre un schéma normalisé.

\vfill
\begin{redbullet}
  \item Identifier la clé primaire de chaque relation (il doit \alert{toujours} y en avoir une)
 \item Identifier les clés étrangères; comprendre quelle clé primaire elles référencent.
 \item Comprendre comment les nuplets sont liés transitivement les uns aux autres.
\end{redbullet}

\vfill

Ces deux bases sont accessibles en ligne: \url{http://deptfod.cnam.fr/bd/tp/sql}
\end{frame}

